\documentclass[11pt,a4paper]{article}
\usepackage{operationsform}
\begin{document}
\OPStitle{科学的不確実性・誇張防止チェックリスト}{TPL-CHK-007}{0.2.0}
\OPSsection{1}{確認項目}
\OPScheck{1}{本文・見出し・図における「観測、推論、仮説、予測、目標」を適切に区分した}
\OPScheck{2}{相関から因果を断定していない}
\OPScheck{3}{研究対象、人数、期間、比較対象と適用範囲を必要な箇所に示した}
\OPScheck{4}{数値の見せ方に効果量や不確実性の重要な欠落がない}
\OPScheck{5}{研究結果から臨床・社会実装までの距離を過度に縮めていない}
\OPScheck{6}{発表の新規性と既存研究の連続性・異なる解釈を検討した}
\OPScheck{7}{タイトル・短文・図表の表現が原稿本文より過度に強くなっていない}
\OPScheck{8}{検証が必要な表現には主張ID、調査担当と修正期限を付した}
\OPSsection{2}{点検結果の記録}
\OPSblankshort{対象案件・対象版}
\OPSblankshort{担当者と点検日}
\OPSblanklong{未解決事項と対応}
\par\Needspace{4\baselineskip}\noindent\textbf{確認に用いた資料ID}\par\vspace{1mm}
{\fontsize{8.5}{13}\selectfont
\begin{tabularx}{\linewidth}{@{}XXXX@{}}
\textbf{対象・項目} & \textbf{内容・参照} & \textbf{担当・状態} & \textbf{備考}\\[2mm]
\strut & \strut & \strut & \strut\\[5mm]
\strut & \strut & \strut & \strut\\[5mm]
\end{tabularx}
}\vspace{2mm}\par
\end{document}
